\section{A shared false length recurrence for three binary scanners}
\label{sec:a289004}

The OEIS entries A289001, A289035, and A289239 give three replacement rules
starting with $w_0=00$ and conjecture that their iterate lengths are the
independently defined sequence A289004
\cite{A289001,A289004,A289035,A289239}.
We refute all three mapping assertions and prove a structural obstruction:
their common actual length sequence satisfies no eventual monic
constant-coefficient integer recurrence, of any order.
The assertion repeated in A289004 and A289001 is counted once.

\subsection{The literal sources and the first disagreement}

Every rule has $00\mapsto0010$ and $10\mapsto010$; the remaining image is
as follows.
\begin{center}
\begin{tabular}{lll}
\hline
Source word & Image of $01$ & Formal variant\\
\hline
A289001 & $001$ & \texttt{left}\\
A289035 & $010$ & \texttt{middle}\\
A289239 & $100$ & \texttt{right}\\
\hline
\end{tabular}
\end{center}
The scanner consumes nonoverlapping matches from the original input, copies
an unmatched leading letter, and retains a final singleton. Replacements
are not rescanned during the same pass. A289035 explicitly specifies the
retained-final-letter convention.

The proposed comparison sequence has
\[
 (a_0,\ldots,a_9)=(2,4,7,11,18,31,52,89,151,257)
\]
and, for $n\ge10$,
\begin{equation}
 a_n=2a_{n-1}-a_{n-2}+2a_{n-4}-a_{n-5}
       +a_{n-7}-a_{n-8}-a_{n-10}.
 \label{eq:a289004:claimed}
\end{equation}
The source prose accidentally repeats index 9 for the last two initial
values. Its displayed terms and Mathematica program unambiguously put 151
at index 8 and 257 at index 9; these are the values used in the formal
independent definition.

\paragraph{The source refutations.}
The three scanners have the same length $L_n$ at every iteration. For
every $n<22$, $L_n=a_n$, while
\[
 L_{22}=268184\ne268182=a_{22}.
\]
Moreover, recurrence~\eqref{eq:a289004:claimed}, when applied to the actual
lengths, fails after every finite cutoff.

The first differing values are:
\begin{center}
\begin{tabular}{r|r|r}
$n$ & actual $L_n$ & A289004 $a_n$\\
\hline
19 & 53908 & 53908\\
20 & 92027 & 92027\\
21 & 157099 & 157099\\
22 & 268184 & 268182\\
23 & 457819 & 457812
\end{tabular}
\end{center}
This is a failure of the proposed recurrence beyond a long common initial
segment, rather than a change in the seed or the starting index.

\subsection{The exact all-iteration count equations}

Call a finite word good if every 1 is immediately followed by a 0.
Equivalently, it is a concatenation of the blocks $0$ and $10$.
Every nonempty good word ends in 0 and has no factor $11$.
The seed is good, and each of the three scanners preserves this property.
For the left variant, the only output block ending in 1 is the image
$001$ of $01$. In a good input, such a pair must be followed by a 0,
and the output of the remaining suffix starts in 0. This checks the only
potential boundary problem; all other relevant output blocks end in 0.
The formal preservation argument treats every good finite input word.

Let $O_n$ denote the number of ones in the $n$th word and put
$e_n=L_n\bmod2$. A good word of length $L_n$ is scanned into exactly
$\lfloor L_n/2\rfloor$ pairs, followed by $e_n$ final zeros.
Every consumed pair produces one 1. A $00$ pair has output length 4,
and a mixed pair has output length 3. Since every input 1 belongs to
a mixed pair, this gives the exact equations
\begin{equation}
 O_{n+1}=\left\lfloor\frac{L_n}{2}\right\rfloor,
 \qquad
 L_{n+1}=2L_n-O_n-e_n,
 \qquad (L_0,O_0)=(2,0).
 \label{eq:a289004:counts}
\end{equation}
The right-hand sides do not depend on the variant. Induction therefore
proves the common length and common ones-count assertions. Eliminating
$O_n$ yields a corrected scalar recurrence:
\begin{equation}
 L_{n+2}+\left\lfloor\frac{L_n}{2}\right\rfloor
       +(L_{n+1}\bmod2)=2L_{n+1},
 \qquad L_0=2,\quad L_1=4.
 \label{eq:a289004:corrected}
\end{equation}
Evaluating this proved count state and the independent recurrence
\eqref{eq:a289004:claimed} establishes both the first discrepancy and
agreement at every earlier index. The Lean proof does not presume
\eqref{eq:a289004:counts}: it derives it from the literal lists.

\subsection{Uniform discrepancy and growth}

Set
\[
 \mu=1-\frac{\sqrt2}{2},\qquad
 \lambda=1+\frac{\sqrt2}{2},\qquad
 S_n=L_n-(2+\sqrt2)O_n.
\]
From \eqref{eq:a289004:counts} and
$2O_{n+1}+e_n=L_n$, direct substitution gives
\begin{equation}
 S_{n+1}=\mu S_n+(1-\mu)e_n.
 \label{eq:a289004:error}
\end{equation}
Since $S_0=2$, $0<\mu<1$, and $0\le e_n\le1$, induction proves
\begin{equation}
 0\le S_n\le2\qquad(n\ge0).
 \label{eq:a289004:bound}
\end{equation}
Every good word has at most half its letters equal to 1. Combining
$2O_n\le L_n$ with \eqref{eq:a289004:counts} gives
\[
 L_{n+1}\ge L_n+\left\lfloor\frac{L_n}{2}\right\rfloor,
 \qquad L_n\ge n+2.
\]
Thus $L_n\to\infty$. A second rearrangement of the exact count equations is
\begin{equation}
 \frac{L_{n+1}}{L_n}
 =\lambda+\mu\frac{S_n}{L_n}-\frac{e_n}{L_n}.
 \label{eq:a289004:ratio}
\end{equation}
The two error quotients tend to zero by \eqref{eq:a289004:bound}, so
\begin{equation}
 \lim_{n\to\infty}\frac{L_{n+1}}{L_n}
 =1+\frac{\sqrt2}{2}.
 \label{eq:a289004:limit}
\end{equation}
This argument proves existence of the limit. For every fixed
$j\ge0$, multiplying $j$ successive ratios proves
$L_{n+j}/L_n\to\lambda^j$.

\subsection{The obstruction to every eventual integer recurrence}

\paragraph{Theorem.}
For every monic polynomial $p(X)=\sum_i p_iX^i\in\mathbb Z[X]$
and every cutoff $N$, there is an $n\ge N$ such that
\[
 \sum_i p_iL_{n+i}\ne0.
\]
In particular, no recurrence
$L_{n+d}=\sum_{i<d}c_iL_{n+i}$ with fixed $c_i\in\mathbb Z$
can hold eventually.

\paragraph{Proof.}
If such a relation held for every sufficiently large $n$, division by
$L_n$ and the fixed-shift limits would give
\[
 p(\lambda)=\sum_i p_i\lambda^i=0.
\]
Monicity would make $\lambda$ an algebraic integer. Closure under
subtraction and multiplication would then make
\[
 (\lambda-1)^2=\frac12
\]
an algebraic integer. The rational-root theorem says that a rational
algebraic integer is an integer, contradicting $1/2\notin\mathbb Z$.
\hfill$\square$

The characteristic polynomial of the precise source recurrence is
\[
 p_0(X)=X^{10}-2X^9+X^8-2X^6+X^5-X^3+X^2+1.
\]
It is monic, so the theorem proves failures of the source recurrence
arbitrarily far out, not merely at iteration 22. The coefficient scope
is exactly monic integer recurrences; arbitrary rational or real
coefficients are not covered by the stated result.

\subsection{Prior results and formal endpoints}

Dekking and Keane proved that the infinite fixed word A289001 is the
Pell word \cite[Theorem~4]{DekkingKeane}. The related position results
are already known. Dekking also proved a floor recurrence for the
final-letter-removed iteration in A293077 \cite{A293077}; the related
density discussion is recorded in A289037 \cite{A289037}.
Those are prior results, not additional new assertions of this work.
The removed-letter lengths differ from the literal retained-singleton
lengths already at iteration 2, where they are 6 rather than 7.

The current source entries, both web-search engines, the trureturing
catalogue and its relevant replacement report, the existing repository,
and global/ProveIt GitHub code searches were checked. No matching earlier
refutation of these retained-singleton length assertions was found in
that bounded screen. Unindexed or differently phrased predecessors
are not excluded.

The seven-module Lean package proves the entire source chain.
\texttt{word\_counts} is the all-index literal scanner bridge;
\texttt{first\_failure} certifies the least failing index;
\texttt{three\_source\_conjectures\_false} states the three mapping
refutations; and \texttt{literal\_corrected\_recurrence} states
\eqref{eq:a289004:corrected}. The analysis endpoints are
\texttt{discrepancy\_bounds} and
\texttt{literal\_length\_ratio\_tendsto}. The stronger source-facing
obstruction is
\begin{center}
\texttt{literal\_no\_eventual\_monic\_relation}.
\end{center}
The exact order-ten specialization is
\begin{center}
\texttt{original\_recurrence\_fails\_after\_every\_cutoff}.
\end{center}
The formalization uses the pinned versions in
Section~\ref{sec:verification}. Its main theorems use
only the standard axioms \texttt{propext}, \texttt{Classical.choice}, and
\texttt{Quot.sound}; no extra axiom or native-decision shortcut is used.
